\documentclass{article}
\usepackage{amsmath}
\usepackage{amssymb}
\usepackage{amsthm}
\usepackage{mathtools}
\usepackage{geometry}
\usepackage{hyperref}

\newcommand{\rk}{\operatorname{rank}}
\newcommand{\GL}{\operatorname{GL}}
\newcommand{\R}{\mathbb{R}}
\newcommand{\Mat}{\mathrm{Mat}}
\newcommand{\Stab}{\operatorname{Stab}}

\theoremstyle{definition}
\newtheorem{theorem}{Theorem}[section]
\newtheorem{proposition}[theorem]{Proposition}
\newtheorem{lemma}[theorem]{Lemma}
\newtheorem{corollary}[theorem]{Corollary}
\newtheorem{definition}[theorem]{Definition}
\newtheorem{remark}[theorem]{Remark}
\newtheorem{conjecture}[theorem]{Conjecture}

\title{First Conjecture Draft}
\date{}

\begin{document}
\maketitle

\section{First Main Conjecture}

This file records successive versions of the first main conjecture in the new direct weight-space setup.

\subsubsection*{Informal version}

The cyclic equations should define a weight-space variety whose orbit decomposition already contains the information needed to stratify the corresponding critical locus. The missing condition that turns a point of the cyclic locus into a genuine critical point should be recoverable by moving along the correct symmetry group.

\subsubsection*{Geometric setup}

Fix:
\begin{itemize}
\item a rank $r$;
\item a subset $S\subset \{1,\dots,r_{\max}\}$ with $|S|=r$;
\item the matrix $W_0=W_S=\Sigma_{XY}P_Q$;
\item the loci $\Gamma_{d,S}^r$ and $\mathcal C_{d,S}^r$ defined in the previous setup files.
\end{itemize}

\subsubsection*{Conjecture A: orbit-surjectivity}

\begin{conjecture}
For every point $W\in \Gamma_{d,S}^r$, there exists an element
\[
g\in G_d^{W_0}
\]
such that
\[
\mu(g\cdot W)=P_SW^\ast.
\]
Equivalently, every $G_d^{W_0}$-orbit in $\Gamma_{d,S}^r$ meets the critical locus $\mathcal C_{d,S}^r$.
\end{conjecture}

\subsubsection*{Conjecture B: orbit-bijection}

\begin{conjecture}
Define
\[
H_S:=\Stab_{G_d^{W_0}}(P_SW^\ast).
\]
Then every $G_d^{W_0}$-orbit in $\Gamma_{d,S}^r$ meets $\mathcal C_{d,S}^r$ in exactly one $H_S$-orbit. Equivalently,
\[
\Gamma_{d,S}^r / G_d^{W_0}
\xrightarrow{\ \sim\ }
\mathcal C_{d,S}^r / H_S.
\]
\end{conjecture}

\begin{remark}
Conjecture B packages two statements:
\begin{itemize}
\item every orbit of $\Gamma_{d,S}^r$ contains a critical point with end-to-end map $P_SW^\ast$;
\item if two points of $\mathcal C_{d,S}^r$ are equivalent through an element of $G_d^{W_0}$ inside $\Gamma_{d,S}^r$, then they differ by an element of the smaller stabilizer group $H_S$.
\end{itemize}
\end{remark}

\subsubsection*{Preferred first main conjecture}

\begin{conjecture}[Orbit correspondence between the cyclic locus and the critical locus]
Fix $r$ and $S$. Let
\[
\Gamma_{d,S}^r
=
\Sigma_d^r\cap \Gamma_{d,S}
\]
be the rank-$r$ weight-space cyclic locus, and let
\[
\mathcal C_{d,S}^r
=
\{W\in \Gamma_{d,S}^r : \mu(W)=P_SW^\ast\}
\]
be the corresponding rank-$r$ critical locus. Then the inclusion
\[
\mathcal C_{d,S}^r \hookrightarrow \Gamma_{d,S}^r
\]
does not induce a quotient by $G_d^{W_0}$ on the critical locus itself, because $\mathcal C_{d,S}^r$ is not stable under that full group when $r>0$. Instead, define
\[
H_S:=\Stab_{G_d^{W_0}}(P_SW^\ast),
\]
where
\[
G_d^{W_0}
=
\{g\in G_d : g_0W_0 g_L^{-1}=W_0\}.
\]
Then every $G_d^{W_0}$-orbit in $\Gamma_{d,S}^r$ meets $\mathcal C_{d,S}^r$ in exactly one $H_S$-orbit. In particular,
\[
\Gamma_{d,S}^r / G_d^{W_0}
\cong
\mathcal C_{d,S}^r / H_S.
\]
Consequently, any stratification of $\Gamma_{d,S}^r$ by $G_d^{W_0}$-stable orbit-type pieces induces a corresponding stratification of $\mathcal C_{d,S}^r$ by intersecting with the critical locus, because each orbit stratum of $\Gamma_{d,S}^r$ contributes exactly one $H_S$-orbit stratum in $\mathcal C_{d,S}^r$.
\end{conjecture}

\subsubsection*{Why this replaces the old lifting statement}

In \texttt{paper\_v0.tex} and \texttt{paper\_v1.tex}, the passage from the cyclic variety to the critical locus went through a residual normal form and a lifting map from residual blocks to weight matrices. The new conjecture avoids that detour:
\begin{itemize}
\item the source and target spaces now both live directly in $\Omega$;
\item the relation between them is expressed purely by an orbit statement in weight space;
\item the projected student map $P_SW^\ast$ replaces the old gauge-fixed residual slice as the anchor condition that singles out critical points inside the larger cyclic locus.
\end{itemize}

\subsubsection*{Sanity check}

The hidden-layer gauge group $G_{\mathrm{hid}}$ cannot be the acting group in this conjecture, because it preserves $\mu(W)$ exactly. Hence it cannot move a general point of $\Gamma_{d,S}^r$ to a point with end-to-end map $P_SW^\ast$ unless that condition already holds.

This is why the first conjecture should be stated using a group with nontrivial endpoint action, such as $G_d^{W_0}$, together with the residual stabilizer subgroup $H_S$ that fixes the distinguished critical product.

\end{document}
